\begin{table*}[t]
\centering
\paperCompactTable
\caption{Comparison with annotation-free GUI agent learning and adaptation methods. We organize the comparison around two limitations addressed by \ourmethod: the lack of a unified mobile interaction and evaluation substrate, and the reliance on isolated rollout experience with coarse or sparse feedback.}
\label{tab:related_work_comparison}
\begin{paperFit}[\textwidth]
\begin{tabular}{l l c c c l c l}
\toprule
\textbf{Method} &
\textbf{Main Source / Setting} &
\multicolumn{3}{c}{\textbf{Pain Point 1: Mobile Interaction and Evaluation Substrate}} &
\multicolumn{3}{c}{\textbf{Pain Point 2: Feedback-to-Optimization}} \\
\cmidrule(lr){3-5} \cmidrule(lr){6-8}
&
&
\makecell{\textbf{Target Mobile}\\\textbf{Interaction}} &
\makecell{\textbf{Auto Task /}\\\textbf{Curriculum}} &
\makecell{\textbf{Hierarchical}\\\textbf{Rollout Eval.}} &
\makecell{\textbf{Feedback}\\\textbf{Granularity}} &
\makecell{\textbf{Cross-Attempt}\\\textbf{Experience}} &
\makecell{\textbf{Policy}\\\textbf{Update}} \\
\midrule
TongUI & Web tutorials & \pmark & \xmark & \xmark & Tutorial-derived trajectories & \xmark & SFT \\
MobileA3gent & User phone trajectories & \pmark & \xmark & \xmark & Auto-annotated instructions & \xmark & SFT / FL \\
GUI-explorer & Autonomous app exploration & \cmark & \pmark & \xmark & Transition-aware knowledge & \pmark & \xmark \\
OS-Genesis & Reverse task synthesis & \pmark & \cmark & \pmark & Trajectory reward model & \xmark & SFT \\
ZeroGUI & Online GUI rollout & \pmark & \cmark & \pmark & Trajectory-level VLM reward & \xmark & Online RL \\
MobileGUI-RL & Online mobile rollout & \cmark & \cmark & \pmark & Trajectory-level composite reward & \xmark & MobiGRPO \\
SEAgent & Desktop software & \xmark & \cmark & \pmark & Step-wise assessment & \pmark & GRPO + imitation \\
ACuRL & Computer-use environments & \xmark & \cmark & \pmark & Outcome evaluator / CUAJudge & \pmark & Curriculum RL \\
UI-Oceanus & Synthetic GUI dynamics & \pmark & \cmark & \xmark & Transition objectives & \xmark & CPT + SFT \\
\midrule
\rowcolor{oursrowcolor}
\textbf{\ourmethod} &
\textbf{Target mobile apps} &
\cmark &
\cmark &
\cmark &
\makecell[l]{\textbf{Trajectory outcome}\\\textbf{+ step-level process}\\\textbf{+ corrective hints}} &
\cmark &
\textbf{\hifpo} \\
\bottomrule
\end{tabular}
\end{paperFit}
\vspace{2pt}
\begin{flushleft}
\footnotesize
\cmark: explicitly supported; \pmark: partially supported or not specific to annotation-free mobile adaptation; \xmark: not supported or not the focus. \mobilegym covers the interaction and evaluation substrate, while \hifpo covers cross-attempt experience reuse and policy optimization.
\end{flushleft}
\end{table*}
